\chapter{Higher order methods I: Runge-Kutta family}

So far in the course we have mostly focused on 
the forward and backward Euler methods both of 
which are first order consistent and convergent. We typically 
just say that these methods are first order accurate to mean 
that their local truncation and global errors are 
$\mathcal{O}(\Delta t)$. While these methods are simple and easy 
to understand, in practice we often would like to use more accurate 
methods, ideally those who have accuracy $\mathcal{O}(\Delta t^p)$ 
for $p=2$ or $p=4$. Let's say we come up with a second order method, 
i.e., $p=2$. Then by halving the stepm size $\Delta t \leftarrow \frac{\Delta t}{2}$, the error should decrease by a factor of $2^2 = 4$. If the method 
is fourth order, i.e., $p=4$ then the error will drop by 
a factor of $16$! Thus, pursuing such methods can realyl pay off in practice. 

In this lecture we will consider one approach to the design of such methods 
that is often called a one-step method: these are 
update schemes where $\ub_{n}$ is computed only from $\ub_{n-1}$. In 
the next example we will introduce multi-step methods where $\ub_{n}$ 
is computed using multiple previous values such as 
$\ub_{n-1}, \ub_{n-2}, \dots, \ub_{n-k}$ for some $k \ge 1$.


\section{From quadrature to Runge-Kutta methods}

Our starting point in deriving higher order methods is 
the approximation of integrals, which is called 
quadrature. Broadly speaking, consider some function 
$g:[a, b] \to \R$ and suppose we want to approximate 
\begin{equation*}
    \int_a^b g(x) \, dx.
\end{equation*}
In quadrature we consider $s$ distinct points $x_1, x_2, \dots, x_s \in [a, b]$ and corresponding weights $w_1, w_2, \dots, w_s$ to approximate the integral as
\begin{equation*}
    \sum_{i=1}^s w_i g(x_i) \approx \int_{a}^{b} g(x) \, dx.
\end{equation*}
The choice of the weights and points $(x_i, w_i)$ constitutes 
the quadrature rule. Popular examples are: 
\begin{equation*}
    \begin{aligned}[t]
        \text{left point rule:} & \quad \int_a^b g(x) \, dx \approx (b-a) g(a), \\
        \text{right point rule:} & \quad \int_a^b g(x) \, dx \approx (b-a) g(b), \\
        \text{midpoint rule:} & \quad \int_a^b g(x) \, dx \approx (b-a) g\Big(\frac{a+b}{2}\Big), \\
        \text{trapezoidal rule:} & \quad \int_a^b g(x) \, dx \approx \frac{b-a}{2} \Big(g(a) + g(b)\Big)
    \end{aligned}
\end{equation*}
These rules are illustrated in \Cref{fig:quadrature-rules}.

\begin{figure}[htp]
    \centering
    \begin{tikzpicture}[thick, scale=0.9,
        declare function={g(\x) = 0.8 + 0.5*\x + 0.4*sin(deg(2*\x)); xa = 0.5; xb = 3; xm = 1.75;},
        area/.style={fill=blue!20, draw=blue!70!black},
        dot/.style={circle, fill, inner sep=1.5pt}]
        % axes, the function g, and labels common to all panels
        \newcommand{\quadpanel}[1]{
            \draw[->] (0,0) -- (4,0) node[right] {$x$};
            \draw[->] (0,0) -- (0,3) node[above] {$y$};
            \draw[very thick, domain=0.1:3.6, samples=100, smooth] plot (\x,{g(\x)}) node[right] {$g(x)$};
            \node[below] at (xa,0) {\strut$a$};
            \node[below] at (xb,0) {\strut$b$};
            \node at (2,-1.1) {#1};
        }

        % left point rule
        \begin{scope}[shift={(0,5)}]
            \draw[area] (xa,0) rectangle (xb,{g(xa)});
            \quadpanel{left point rule}
            \node[dot] at (xa,{g(xa)}) {};
        \end{scope}

        % right point rule
        \begin{scope}[shift={(6,5)}]
            \draw[area] (xa,0) rectangle (xb,{g(xb)});
            \quadpanel{right point rule}
            \node[dot] at (xb,{g(xb)}) {};
        \end{scope}

        % midpoint rule
        \begin{scope}[shift={(0,0)}]
            \draw[area] (xa,0) rectangle (xb,{g(xm)});
            \draw[dashed, thin] (xm,0) -- (xm,{g(xm)});
            \quadpanel{midpoint rule}
            \node[below] at (xm,0) {$\frac{a+b}{2}$};
            \node[dot] at (xm,{g(xm)}) {};
        \end{scope}

        % trapezoidal rule
        \begin{scope}[shift={(6,0)}]
            \draw[area] (xa,0) -- (xa,{g(xa)}) -- (xb,{g(xb)}) -- (xb,0) -- cycle;
            \quadpanel{trapezoidal rule}
            \node[dot] at (xa,{g(xa)}) {};
            \node[dot] at (xb,{g(xb)}) {};
        \end{scope}
    \end{tikzpicture}
    \caption{The left point, right point, midpoint, and trapezoidal rules 
    applied to a smooth function $g$ on $[a, b]$. The shaded region in 
    each panel is the area computed by the quadrature rule, while the 
    exact integral is the area under the curve $g(x)$. Black dots mark 
    the points where $g$ is evaluated.}
    \label{fig:quadrature-rules}
\end{figure}

Let us now consider a generic scalar valued ODE 
\begin{equation*}
    u' = f(t,u), 
\end{equation*}
where we switched to the standard prime notation 
$u'(t) = \frac{d u}{dt}(t)$ for convenience.
We will develop our methods for this case but note that 
they have a natural extension to systems of ODEs that 
we will discuss later. 

Consider the interval $[t_{n-1}, t_n]$ and integrate both sides of the ODE:
\begin{equation*}
    u(t_n) - u(t_{n-1}) = 
    \int_{t_{n-1}}^{t_n} u'(t) \, dt 
\end{equation*}
Let us now approximate the integral on the 
right hand side using the left point rule: 
\begin{equation*}
    u(t_n) - u(t_{n-1}) \approx (t_n - t_{n-1}) u'(t_{n-1}) 
    = \Delta t \, f(t_{n-1}, u(t_{n-1})).
\end{equation*}
We recognize this as the forward Euler scheme. Using the 
right point rule we obtain backward Euler:
\begin{equation*}
    u(t_n) - u(t_{n-1}) \approx \Delta t u'(t_n) 
    = \Delta t f(t_n, u(t_n)) 
\end{equation*}
Continuing forward, if we use the midpoint rule we 
obtain a new method, which is simply known as the (implicit)
 midpoint method: 
\begin{equation*}
    \begin{aligned}
    u(t_n) - u(t_{n-1}) &\approx 
    \Delta t u'\Big(\frac{t_{n-1}+t_n}{2}\Big) 
    = \Delta t f\Big(\frac{t_{n-1}+t_n}{2}, u\Big(\frac{t_{n-1}+t_n}{2}\Big)\Big)\\
    &\approx \Delta t f \Big( t_{n- \frac12}, \frac{u(t_n) + u(t_{n-1})}{2} \Big).
    \end{aligned}
\end{equation*}


% \begin{equation*}
%     u(t_n) - u(t_{n-1}) = \int_{t_{n-1}}^{t_n} u'(t) \, dt = \int_{t_{n-1}}^{t_n} f(u(t)) \, dt.
% \end{equation*}



\bibliographystyle{plainnat}
\bibliography{references}